\section{Design-based abstention and a shared population error budget}
\label{sec:design-budget}

The version-8 source-failure branch used only the target outcome allowance
$\alpha_T=.01$, although no sources were used after failure. This needlessly
widened the fallback relative to the separately reported target-only interval.
The correction must account for the distinction between conditional outcome
events and the unconditional target-composition event.

\subsection{Why independently valid branch intervals are insufficient}

Let $B(\mathcal D)$ be the design-only decision to borrow. Both the valid-count
screen and balance-weight screening can depend on the observed target
covariates through $b_t$. A composition concentration event is unconditional
over these covariates. Thus switching between two separately calibrated
population intervals using $B(\mathcal D)$ does not automatically preserve
their individual coverage level. In particular, increasing the composition
allowance only on the fallback branch cannot be justified by conditioning on
the same design.

Instead fix a single coefficient allowance $\alpha_R$ and a single composition
allowance $\alpha_X$, common to all branches. Put
$\alpha_C=\alpha-\alpha_R-\alpha_X$. Conditional on $\mathcal D$, split
$\alpha_C$ between source noise, target noise and dispersion when borrowing,
and assign all of it to target noise when abstaining. Both branches then
satisfy
\[
 \Prb\{\tau_X\notin I_X(B(\mathcal D))\mid\mathcal D\}\le\alpha_C.
\]
One common coefficient/composition event completes the population argument:
\[
 \Prb\{\tau_t\notin I\}\le\alpha_C+\alpha_R+\alpha_X=\alpha.
\]
No independence of the decision and target covariates is required. Estimated
approximation certificates still need their separate failure allowance.

The new \texttt{shared\_composition} allocation uses
\[
 \alpha_R=.005,\qquad\alpha_X=.0225,\qquad\alpha_C=.0225.
\]
On the borrowing branch, preserve the old conditional proportions:
\[
 (\alpha_S,\alpha_T,\alpha_D)
   =\frac{.0225}{.035}(.01,.01,.015).
\]
On the abstention branch, use $\alpha_T=.0225$ alone. Its reported interval
is then \emph{exactly} the complete-budget protected target-only interval,
including the same composition certificate.

The legacy allocation remains an explicit control. A second control,
\texttt{conditional\_reallocation}, keeps its common $\alpha_X=.01$ and
reassigns the full $.035$ conditional budget to target noise only on design
fallback. It is valid but generally does not equal the previous standalone
target interval, whose composition allowance is $.0225$.

This change guarantees equality with target-only on pre-outcome abstention.
It does not guarantee that every realized borrowing interval is shorter than
the complete-budget target interval. Selecting the shorter of differently
calibrated population procedures after seeing outcomes would require a new
joint coverage argument.

\subsection{Precision-aware rules using design information only}

Numerical estimability does not ensure a useful source estimate. For arm
$a$ at site $j$, define weight inflation
\[
 \mathcal I_{ja}=n_{ja}\|\gamma_{ja}\|_2^2\ge1.
\]
It equals one for uniform arm weights. A prespecified threshold $L_I$ removes
a site if either arm exceeds it. The experiments declare $L_I=4$ as a main
diagnostic and report $2,8$ sensitivities, without selecting thresholds by
observed coverage or length. The count update still uses the original
guarantee: $g_a'=\max(0,g_a-m)$ after all design exclusions.

A separate optional precision gate uses the exact retained design factors
but sets the \emph{observed} source dispersion to zero in the certificate.
Let $B_{\rm ref}(\mathcal D)$ be its resulting bias radius and let
$r_S(\mathcal D)$ be the borrowing-branch noise radius. Abstain if
\[
 r_S(\mathcal D)+B_{\rm ref}(\mathcal D)
 \ge \sqrt{\frac{\log(2/\alpha_C)}2\sum_a S_{ta}}.
\]
Approximation modes include their design-based centering/MGF budgets when
forming $B_{\rm ref}$. This reference calculation is a heuristic for likely
precision under small observed dispersion. It is not a deterministic lower
bound on the realized source radius, which also depends on outcomes. Its
coverage validity follows solely from being fixed by design before outcomes;
it can discard a borrowing opportunity. The full-target fallback therefore
has a clear interpretation without an unsupported claim of optimal screening.

\subsection{An observable sufficient condition for the conditional rate}

The weight guard has a theoretical role beyond numerical stability. Suppose
all retained arms have at least $n_*$ patients, inflation at most $L_I<\infty$
and leverage at most $h_*<1$, with constants fixed along the sequence. Then
\[
 S_j\le\frac{L_I}{n_*},\qquad
 d_j^{\max}\le\frac{L_I}{(1-h_*)n_*},\qquad
 G_j,F_j\le\frac{L_I^2}{(1-h_*)^2n_*^2}.
\]
Indeed, $\max_i\gamma_i^2\le\sum_i\gamma_i^2=S_j$ and
$\sum_i\gamma_i^4\le S_j^2$. These bounds are weaker than the
$n_s^{-3}$ and $n_s^{-2}$ sufficient bounds used in the main theorem, but
still imply for $K'$ retained sources
\[
 L=O\{(n_*K')^{-1}\},\quad
 Q_0=O\{n_*^{-2}(K')^{-1}\},\quad
 c=O\{(n_*K')^{-1}\}.
\]
Under exact linear means, positive adjusted valid proportions and arm
$D=O\{n_*^{-1}(K')^{-1/2}\}$, the same argument gives
\[
 \E(|J_S|\mid\mathcal D)=O\{n_*^{-1/2}(K')^{-1/4}\}.
\]
This is a directly checkable sufficient
condition on the retained design, without requiring an individual-weight
bound of order $n_*^{-1}$. Dimension can still affect the probability that
enough sites pass, the minimum retained arm size and the target-composition
cost. Fixed thresholds do not remove these issues, and do not guarantee useful
constants in a small-sample experiment.
